The case study examines the domestic airline industry in Brazil from 2012 to
2016. GOL Linhas A\'{e}reas Inteligentes (GOL) and LATAM Airlines (LATAM)
were the dominant airlines throughout this period.\footnote{%
Effective June 2012, Brazilian Airline TAM Linhas A\'{e}reas purchased
Chilean LAN Airlines to form the LATAM Airlines Group. Thus, there will
sometimes be reference to LAN Airlines and TAM.} The other two key airlines
are Aerov\'{\i}as del Continente Americano (Avianca) and Azul Linhas
A\'{e}reas Brasileiras (Azul) who, while starting small in 2011, saw their
market shares grow significantly. Figure~\ref{fig:airline_capacities} reports the four airlines'
capacities over 2011-16.

\begin{figure}[htbp]
\centering
\includegraphics[width=\textwidth]{../data/outputs/figures/ask_by_airline_2011_2016_16x9.pdf}
\caption{Airline capacity in Brazil, 2011--16}
\label{fig:airline_capacities}
\begin{minipage}{\textwidth}
\vspace{1em}
\footnotesize
\textit{Notes:} Available seat kilometers (ASK) for GOL, LATAM, Avianca, and Azul on
domestic regular flights, aggregated to airline-year totals from ANAC data \cite{anac2011_2016}.
\end{minipage}
\end{figure}

The time period 2012--16 was selected based on the LLM's output. Over the
entire database spanning from 2002 to 2025, there are 25 LLM-flagged
GOL transcripts (score of at least 75) and 20 of those occurred in 2012--16 with
only one over 2017--21 and four during 2022--23.\footnote{%
LLM-flagged transcripts for GOL: 2012 (4 transcripts), 2013 (7), 2014 (4),
2015 (2), 2016 (3), 2017 (0), 2018 (1), 2019 (0), 2020 (0), 2021 (0), 2022
(3), 2023 (1).} The LLM-validated step screened 15 GOL transcripts (with a
mean score of at least 75 based on at least 10 runs) and all but one of them
occurred during 2012--16.\footnote{%
LLM-validated transcripts for GOL: 2012 (4 transcripts), 2013 (5), 2014 (1),
2015 (2), 2016 (2), 2022 (1).} For LATAM, it had eight LLM-flagged
transcripts with six of them over 2012--16 (of which three are LLM-validated)
and two in 2017 (of which one is LLM-validated). In sum, the LLM\ is
identifying a large number of suspicious announcements during 2012--16, none
in the preceding years, and very few in the immediately ensuing years.

We limited our analysis to the 39 earnings calls for GOL\ and LATAM over
2012 to 2016.\footnote{%
For purposes of brevity, we are excluding transcripts that are not earnings
calls. Research assistant Vivian Zhang read the entirety of the 39 earnings
calls and collected excerpts with possible collusive content. Co-author
Joseph Harrington then read the excerpts upon which the analysis of this
section is based.} During that time period, there were no earnings calls for
Avianca and Azul.\footnote{%
Due to when Avianca was listed on the NYSE\ and various restructuring, the
database has 11 earnings calls from 2017 Q2 to 2019 Q4 and nine earnings
calls from 2023 Q1 to 2025 Q1. Azul was not listed on the NYSE\ until 2017
and the database has 33 earnings calls from 2017 Q1 to 2025 Q1.} Provided
below in chronological order are all of the earnings calls with any relevant
content. When the content is conveying a similar message, they are placed
together under a bullet point with a one or two sentence description of that
message. For each earnings call, a succinct description of the content of
the earnings call is provided along with the date (in the format
year/month/day). For example, GOL earnings calls for March 27, 2012 and May
4, 2012 both convey that GOL is limiting its supply. So that the reader can
draw their own interpretations, relevant excerpts from the LLM-identified
passages are provided in the footnotes.\footnote{%
To ease reading, many excerpts have been condensed and combined but always
in a manner not to affect their meaning. In cases where the meaning of the
content is unclear, the original passage was reproduced unaltered.}

\begin{itemize}
\item GOL states it will limit supply.

\begin{itemize}
\item 2012/03/27 -- GOL: We are committed to responsibly adding supply.
\footnote{%
``The company has already affirmed its commitment to the responsible
addition of supply on the domestic market and recently announced a 0\%
supply growth. We are therefore announcing a capacity cut as of March.''}

\item 2012/05/04 -- GOL: We are reducing supply and the industry will
restore profitability by limiting supply.\footnote{%
``GOL's new target of reducing domestic supply by up to 2\% over 2012,
together with the industry adoption of a conservative approach to increase
supply, are indicative of a market governed by rationality in order to
resume profitability. The long-term winners are the ones that have the
courage of not saying that they are buying a lot of planes in the short
run.''}
\end{itemize}

\item GOL and LATAM express the industry needs to cut capacity and note that
they are cutting capacity but some firms are not cutting capacity.

\begin{itemize}
\item 2012/08/13 -- LATAM: The market has suffered from overcapacity but
conditions are in place for the industry to limit capacity. We will decrease
capacity.\footnote{%
``Brazil has historically suffered from overcapacity. However, we believe
that conditions are in place for carriers to sustain capacity ceiling and
focus increasingly on profitability. We currently expect decreased capacity
between 2\% and 3\% this year.''}

\item 2012/08/14 -- GOL: We are increasing our capacity reduction and we
expect competitors to be less aggressive in pricing. But our market share
has been declining because some firms are expanding.\footnote{%
``Previously, we're saying that the decrease in capacity would be 2\% for
the year. And now, we are saying that the new range is between 2\% and 4.5\%
on the negative side. We do expect to see our competitors being less
aggressive in their price policies for the following months. But when we see
the total market, it's also clear that our market share in combination with
our main competitor, is going down because the newcomers are investing a lot
in increasing their capacities.''}

\item 2012/11/13 -- LATAM: Some firms are pursuing capacity discipline and
we support it for the industry.\footnote{%
``Some of the major players of the industry are taking a strategy of
capacity discipline that will help to increase the load factor. We're going
to listen to capacity discipline of the industry.''}

\item 2012/11/14 -- GOL: We and a competitor have cut capacity and are
showing capacity discipline, and we intend to cut more.\footnote{%
``We cut the capacity on April in comparison to March by 8\%, which is going
to give us an annual effect of around 4.5\%. The competitor just now is
taking the same measure. We do believe that the current industry discipline
should reduce capacity and we would like to find additional opportunities to
reduce capacity.''}
\end{itemize}

\item GOL and LATAM emphasize that capacity discipline will restore
profitability, and GOL notes that other firms are now following its lead to
reduce capacity.

\begin{itemize}
\item 2013/03/20 -- LATAM: We believe capacity discipline is the key to
restoring profitability.\footnote{%
``We remain convinced that capacity discipline and adequate segmentation of
the market will provide the basis for continued healthy load factors and
significant improvement in operating results in 2013. Our strategy is very
focused on capacity discipline, and we believe that this is key for
regaining profitability for the year in Brazil.''}

\item 2013/03/26 -- GOL: Other firms are now following our lead to reduce
capacity. Less capacity is the path to higher margins.\footnote{%
``For 2013, the company has already announced the curbs of domestic supply
of around 7\%, meaning a reduction of more than 10\% in 2 years, reinforcing
the company's commitment to resuming operating margins. The leading market
players have begun to follow the supply reduction initiated by GOL in April
2012, an important sign that Brazil's airline industry is on the way to
constructing a more sustainable environment.''}
\end{itemize}

\item GOL and LATAM reiterate that reducing capacity is the key to better
market outcomes---such as higher revenue---and they are continuing to
practice capacity discipline.

\begin{itemize}
\item 2013/05/15 -- LATAM: We believe capacity discipline is the key to
higher profits.\footnote{%
``We remain convinced that capacity discipline and adequate segmentation of
the market will provide the basis for continued healthy load factors and
expect continued improvement in operating results in 2013.''}

\item 2013/08/21 -- LATAM: We continue to focus on reducing capacity.
\footnote{%
``We continue focused on reducing capacity. We expect to reduce capacity by
between 7\% and 9\%.''}

\item 2013/11/13 -- GOL: Supply is down and that is key to improving market
outcomes.\footnote{%
``Domestic supply is down year-over-year by 7\%, almost 10\% in the 9 months
of 2013. We think this is absolutely necessary for the market to reorganize
itself.''}

\item 2014/03/18 -- LATAM: Maintaining capacity discipline has led to higher
profits.\footnote{%
``TAM continues to make significant improvement in the financial result of
the domestic Brazil passenger operations, maintaining capacity discipline,
6\% reduction quarter-over-quarter.''}

\item 2014/03/26 -- GOL: The priority is raising PRASK (passenger revenue
per available seat kilometer) and a main driver is reducing capacity and
rationalizing supply. We do not see future industry capacity additions.
\footnote{%
``Our strategy of rationalizing supply is in place and is one of the main
drivers that is causing our revenues to really move up. The priority is to
increase the PRASK at the level of 2 digits. So that strategy has been quite
successful, and it has also helped the markets to drive towards better
capacity discipline. We do not see in the horizon, the time when the
Brazilian domestic market will be able to accept additional capacity
again.''}
\end{itemize}

\item GOL and LATAM say they are staying the course of capacity discipline,
and LATAM notes that others are expected to do so, too.

\begin{itemize}
\item 2014/05/14 -- LATAM: The competitive environment remains intense and
we continue to support a rational capacity strategy.\footnote{%
``We do expect to see continued reductions in ASKs during the second
quarter. The competitive environment in Brazil remains very intense. So we
continue with very rational capacity strategy in the domestic Brazil with a
very positive effect on load factors.''}

\item 2014/05/15 -- GOL: We are rationalizing supply to pursue price
increases.\footnote{%
``We can see GOL's strategy to pursue an increase in price and
rationalization of supply which is currently 12\% lower than when the
company started to cut capacity in May 2012.''}

\item 2014/08/13 -- LATAM: We reduced capacity and expect others to do so,
too.\footnote{%
``What we observed is that there's a lot more capacity additions than
demand. We reduced this capacity, and we expect others to reduce.''}

\item 2014/08/14 -- GOL: The industry is pursuing limitations on capacity.
\footnote{%
``The industry is still moving towards a very strict control on capacity.
The indications we got from the market related to the competitors and to the
industry as a whole doesn't show any important movement to add capacity.
Therefore, I think the (pricing) discipline will continuously offer us a
healthier environment to recover margins.''}

\item 2014/11/12 -- GOL: We will continue to be rational in terms of
capacity.\footnote{%
``We will still be very, very rational in terms of capacity in Brazil. The
strategy has been working for us.''}
\end{itemize}

\item GOL suggests that Avianca and Azul will engage in capacity discipline.

\begin{itemize}
\item 2015/03/31 -- GOL: We believe Avianca and Azul will pursue the
rational sales approach started by GOL and followed by LATAM.\footnote{%
``We are maintaining our firm commitment to not expanding domestic supply in
2015. But I could say that we do see a more rational approach in the future
sales eventually coming from the market as a whole. We do expect that our
competitors will follow this rational approach already started by GOL,
followed by TAM. And we believe that it should be followed by Avianca and
Azul along the following months. We say our roof (for capacity) is 0 growth.
It could be lower. This is basically our clear signage to the market that
we won't increase capacity.''}

\item 2015/05/13 -- GOL: We continue to be committed not to expand supply.
While some competitors have increased supply, we believe this is about to
change.\footnote{%
``We maintain our commitment of not expanding the supply in the domestic
market. There is an overcapacity because some of our competitors have
increased their capacity over the last period, while the demand came down. I
believe that this behavior is about to change when we look at the future
sales -- the inventory. I would say all the competitors have changed the
number of seats being pushed available for sales along the following months.
And this is a quite good signal that this capacity discipline will be
considered more seriously by everyone else.''}
\end{itemize}

\item GOL states that some firms have not followed their lead to reduce
capacity. Both GOL and LATAM say they will continue with capacity cuts.

\begin{itemize}
\item 2015/08/14 -- GOL: We led in cutting capacity but some competitors did
not and there is still excess capacity. We will continue to cut.\footnote{%
``We still do have an over-capacity in the market. Therefore, we are again
leading the market to further capacity decrease. We will reduce supply to
between 2\% and 4\% in the second half. We have started this capacity
reduction already in April 2012. And we have cut, since then, 14\%. LATAM is
somehow catching up the same capacity reduction.''}

\item 2015/08/14 -- LATAM: We are cutting capacity more than originally
planned.\footnote{%
``We were previously expecting to be flat for 2015 in terms of total
capacity and our current estimate is to have a reduction of between 2\% and
4\%.''}
\end{itemize}

\item GOL expresses that other firms seem to be ready to reduce capacity,
and LATAM notes that more capacity needs to be cut.

\begin{itemize}
\item 2015/11/12 -- GOL: The industry is now showing signs of support for
capacity reductions.\footnote{%
``The industry is giving clear signals that there are signs that there is a
different behavior related to capacity. The major difference in comparison
to the first capacity reduction cycle is that now the pressure on the demand
does not allow any capacity increase even from the other competitors.''}

\item 2015/11/13 -- LATAM: We were the only firm to reduce capacity in the
last quarter. More capacity needs to be cut.\footnote{%
``TAM was the only company to reduce capacity in the third quarter of this
year. We see a complicated economic year for Brazil. That's why we're
forecasting another reduction for next year as necessary. Our estimate is to
reduce capacity between 9\% and 6\% for next year.''}
\end{itemize}

\item GOL and LATAM continue to support capacity reductions.

\begin{itemize}
\item 2016/03/09 -- LATAM: We have been leading the industry in capacity
cuts and capacity reduction is occurring throughout the industry.\footnote{%
``And we have been leading the industry in terms of capacity reduction. We
had originally talked about a reduction of between 6\% and 9\%. Now we have
increased that to a reduction of between 8\% and 10\%. What we do see is the
capacity reduction overall in the industry, which we think will have
positive impacts in our yields in the future.''}

\item 2016/03/30 -- GOL: Our supply cuts have been rational for the
industry.\footnote{%
``Together, we are responsible for approximately 14 billion ASKs less in the
system that has been bringing more rationality for the industry.''}
\end{itemize}

\item GOL and LATAM state that they have cut capacity but more needs to be
cut. GOL notes some firms have still not cut capacity.

\begin{itemize}
\item 2016/05/12 -- GOL: In response to overcapacity, we have reduced
capacity but small firms have not. There needs to be more capacity cuts.
\footnote{%
``Overcapacity in Brazil has added pressure on GOL. While the company has
been very disciplined, others have not. GOL has been reducing capacity since
2012. However, smaller players have been adding capacity, negatively
affecting sector results. The discipline regarding capacity is much more
rational than before, but I wouldn't say that it has achieved its balance in
comparison to the demand when we are talking on the whole system. It's
improving but we are still not close to the optimum point.''}

\item 2016/05/12 -- LATAM: We are increasing our capacity reductions but
more industry capacity needs to be cut.\footnote{%
``We are increasing our capacity and we now expect capacity to be reduced
between 10\% and 12\%. We're convinced that this rational approach to
capacity is key to being able to improve our results in that market. Despite
the strong capacity cuts that the market implemented, it was still not
enough to compensate for the demand drop.''}
\end{itemize}

\item GOL expresses how it has led in reducing capacity, some firms have not
followed, and all firms benefit from reduced industry capacity. LATAM
conveys that there is too much industry capacity and it will not add
capacity.

\begin{itemize}
\item 2016/08/12 -- LATAM: We do not expect to raise capacity.\footnote{%
``We continue to expect our passenger capacity to be relatively flat this
year with respect to 2016.''}

\item 2016/08/16 -- GOL: We have led in reducing capacity but two firms have
not followed that strategy. We will continue to cut capacity and believe
lower industry capacity will benefit all firms.\footnote{%
``Although GOL has provided leadership in reducing capacity, industry
profitability has been negatively impacted by the capacity growth of 2
players, one of which has not reduced capacity, and the other has even added
capacity during the downturn in the market. The competitive environment is
still tough. We also believe that the reduction in the industry's capacity
will benefit us all.''}

\item 2016/11/07 -- GOL: We have led in reducing capacity but two firms have
not followed that strategy. We will continue to cut capacity and believe
lower industry capacity will benefit all firms.\footnote{%
``Since 2011, GOL has worked to lead the industry in rational and profitable
capacity growth. In the adjustment of industry capacity to match demand, not
all competitors have cut capacity and reduced frequencies on destinations
and routes in Brazil. Some competitors have continued to cut capacity and
reduced frequencies. Other competitors have continued to add capacity, in
spite of their higher operating costs. And this has had a big impact on
sector profitability. To the extent that everybody is disciplined on
capacity, we should see a solid yield environment.''}

\item 2016/11/11 -- LATAM: There is still too much capacity and we will not
add capacity in the near term.\footnote{%
``We still believe that there is some excess capacity in the market. We're
not thinking about adding any capacity in the next few months.''}
\end{itemize}
\end{itemize}

\medskip

During public announcements spanning from 2012 to 2016, GOL and LATAM
explicitly and repeatedly stated there is too much industry capacity and
reassured each other that they will reduce capacity. Good to their word,
both airlines significantly reduced their capacities as can be seen in Figure~\ref{fig:airline_capacities}. GOL
provided confirmation in their 2016/03/30 earnings call. 

\begin{quote}
``Along the last 4 years ... GOL reduced approximately 14\% of its supply in the domestic markets 
and was followed by its main competitor. Together, we are responsible for approximately 
14 billion ASKs less in the system, a phenomena that has been bringing more rationality 
for the industry. However, it's not the same effect produced by other competitors in our industry ...''
\end{quote}

\noindent As just noted by GOL, Avianca and
Azul did not reduce their capacities and GOL\ and LATAM criticized those two
airlines for not doing their part. Perhaps to encourage their cooperation,
GOL conveyed the expectation that Avianca and Azul were going to get on
board, though apparently they never did. The explicitness of the content and
the repeated back-and-forth between GOL and LATAM is compelling evidence
that they both invited coordinated capacity cuts and both accepted the other
firm's invitation. In other words, they had an agreement to reduce their
capacities which, based on this evidence, is likely to be an antitrust law
violation in many jurisdictions.\footnote{%
Whether the capacities of GOL\ and LATAM were lower because of these
announcements is a separate question about economic effect rather than
liability. To say something about the counterfactual capacities, one would
need to engage in an analysis such as in \textcite{aryal2022coordinated}.
They used variation in the content of earnings calls across airlines and
time to test whether all (legacy) airlines on a route having ``capacity
discipline'' or similar content in their earnings calls in a quarter was
correlated with lower capacity for that route in the ensuing quarter. We do
not have that form of variation in the case of the Brazilian air travel
market.}

The case study illustrates the potential value of screening public
announcements using an LLM. The LLM's screen identified public announcements
for GOL\ and LATAM containing collusive content. This finding was the basis
for conducting a human review of public announcements by GOL and LATAM
during the time period identified by the LLM. That human review delivered
compelling evidence of communications facilitating an illegal agreement. If
it had been found at the time, we believe this evidence would have justified
prosecution of GOL\ and LATAM by the Brazilian competition agency Conselho
Administrativo de Defesa Econ\^{o}mica (CADE).

Interestingly, inspection of earnings calls from 2022--23 shows a resurgence of
disconcerting content. There were four LLM-flagged GOL transcripts and three
of them had relevant content which we provide here.\footnote{%
The 2022/07/28 earnings call passes the LLM-validated screen, while the
2022/03/14 earnings call is just below the threshold. The 2023/07/27
earnings call falls well below the threshold. None of LATAM's transcripts
were flagged.}

\begin{itemize}
\item 2022/03/14 (mean score = 74.17): ``GOL have led the market towards
rationality. The industry has demonstrated an even higher level of
rationality in comparison to what we have seen so far.''

\item 2022/07/28 (mean score = 76.25): ``We're being very disciplined on
capacity given what's going on with costs and particularly fuel costs. I can
say it was an industry cut because we also put 10\%. We saw also competition
cutting 10\% to sustain all the fare activity that we are required to do in
this very high cost environment. What we are seeing in the market is really
a rational behavior in terms of fares.''

\item 2023/07/27 (mean score = 65): ``We want to create the right equilibrium
and demand and supply. We are positioning ourselves as kind of a leading the
capacity discipline, we are cutting. We are seeing that there is rationality
in the market.''
\end{itemize}

\noindent CADE may want to take a close look at the domestic airline
industry.
